\documentclass[a4paper,11pt]{article}
\usepackage[latin1]{inputenc}
\usepackage[italian]{babel}
\usepackage{amsmath}
\usepackage{amsfonts}
\usepackage{amsthm}
\usepackage{listings}
\usepackage{float}
\usepackage{graphicx}

\newcommand{\nats}{\mathbb{N}}
\newcommand{\reals}{\mathbb{R}}
\newcommand{\codegeneric}[1]{``\lstset{basicstyle=\ttfamily}\lstinline@#1@''}
\newcommand{\codexml}[1]{``\lstset{basicstyle=\ttfamily, language=Xml}\lstinline@#1@''}
\newcommand{\codeprolog}[1]{``\lstset{basicstyle=\ttfamily, language=Prolog}\lstinline@#1@''}
\newcommand{\codexmlnb}[1]{\lstset{basicstyle=\ttfamily, language=Xml}\lstinline@#1@}
\newcommand{\codeprolognb}[1]{\lstset{basicstyle=\ttfamily, language=Prolog}\lstinline@#1@}
\newcommand{\filecodeprolog}[1]{\lstinputlisting[basicstyle=\ttfamily, breaklines=true, breakatwhitespace=false, numbers=left, numberstyle=\tiny, frame=single, language=Prolog]{#1}}

\lstnewenvironment{blockcodexml}{\lstset{basicstyle=\ttfamily, breaklines=true, breakatwhitespace=false, numbers=left, numberstyle=\tiny, frame=single, language=Xml}}{}

\newfloat{code}{!h}{}
\floatname{code}{Codice}

\title{
  \includegraphics[width=5cm]{logo-unifi.jpg}\\[1cm]
  \textbf{PARADIGMI DI PROGRAMMAZIONE\\relazione elaborato}
}
\author{\textbf{STEFANO MARTINA}}

\begin{document}
\maketitle
\newpage
\section{Introduzione}
Il programma consiste in un validatore di files \emph{XML} rispetto a certi limiti sintattici che specificher� pi� avanti. Il programma inoltre � in grado di completare un \emph{XML} incompleto sotto certi limiti dovuti al non utilizzo della \emph{ricorsione tail}.

\`E possibile fornire come input al programma, per le funzionalit� di validazione e di completamento, sia un file \emph{XML} che una stringa contente direttamente il codice

Per lo sviluppo ho usato l'implementazione \emph{SWI-Prolog} in ambiente \emph{Gnu Linux} e usando come editor \emph{Emacs}. Il risultato dell'elaborato consiste in una \emph{grammatica formale} anche se sono arrivato alla stesura tramite un approccio sperimentale e non dopo una formalizzazione matematica preventiva.
\section{XML, codifica e limiti}
L'\emph{XML} � un \emph{metalinguaggio} di \emph{markup} che definisce una sintassi per esprimere una struttura ad albero dove ogni nodo si definise \emph{tag}. La caratteristica principale � che � possibile definire una struttura personalizzata a seconda dell'utilizzo che se ne vuole fare, cosa che ne permette un uso molto eterogeneo: dalle pagine web (con l'\emph{xhtml} che � una formalizzazione in \emph{xml} dell'\emph{html}) allo scambio di dati o alla definizione di linguaggi.

La codifica dei caratteri � \emph{unicode} valida in certi range, inoltre ogni carattere � possibile inserirlo con la sua corrispondente codifica decimale \codexml{\&\#...;} o esadecimale \codexml{\&\#x...;}. Ci sono dei caratteri speciali che non � possibile usare nel contenuto, ma solo nella struttura, per inserirli nel contenuto � necessario usare le entit� xml illustrate nella seguente tabella \ref{tab:entita}.
\begin{table}[!h]
  \begin{center}
    \begin{tabular}{|c|c|}
      \hline Carattere & Codifica\\
      \hline \textless & \codexmlnb{\&lt;}\\
      \textgreater & \codexmlnb{\&gt;}\\
      \& & \codexmlnb{\&amp;}\\
      ` & \codexmlnb{\&apos;}\\
      `` & \codexmlnb{\&quot;}\\
      \hline
    \end{tabular}
    \caption{Entit� \emph{xml}}\label{tab:entita}
  \end{center}
\end{table}

Come gi� detto un codice \emph{xml} � costituito da un albero di tags con un singolo tag come radice.
\begin{code}[!h]
  \begin{blockcodexml}
    <nome attributo1="valore1" attributo2="valore2" attributoX="valoreX">
      contenuto
    </nome>
  \end{blockcodexml}
  \caption{Tag generico}\label{cod:singolotag}
\end{code}
Ogni singolo tag � costituito come nel codice \ref{cod:singolotag}, dove \codexml{contenuto} pu� essere vuoto, un testo oppure una serie di altri \emph{tags}, \codexml{nome} e \codexml{attributoX} non devono contenere i caratteri: \\\codegeneric{!"#\$\%&'()*+,/;<=>?\@[\\]^`\{|\}~}, spazi, o cominciare per un numero o per \codegeneric{-} o \codegeneric{.}. Ovviamente vi possono essere un numero qualsiasi di attributi, anche nessuno. Nel caso in cui \codexml{contenuto} sia vuoto, � possibile usare una forma abbreviata come nel codice \ref{cod:abbreviato}.
\begin{code}[!h]
  \begin{blockcodexml}
    <nome attributo1="valore1" attributo2="valore2" attributoX="valoreX" />
  \end{blockcodexml}
  \caption{Tag vuoto abbreviato}\label{cod:abbreviato}
\end{code}

Oltre alla struttura dei \emph{tags}, nello standard pu� essere presente anche una riga preliminare di dichiarazione definita come nel codice \ref{cod:dichiarazione}
\begin{code}[!h]
  \begin{blockcodexml}
    <?xml version="1.0" encoding="UTF-8" ?>
  \end{blockcodexml}
  \caption{Dichiarazione}\label{cod:dichiarazione}
\end{code}
che deve essere necessariamente la prima e dove vi sono le informazioni sulla versione usata dello standard e sulla codifica dei caratteri.

Se un documento \emph{xml} rispetta le semplici regole appena definite si dice che esso � \emph{ben formato}, oltre a questo, � possibile definire uno schema che indica la struttura e i nomi possibili dei tags per esempio come avviene nelle pagine web in \emph{xhtml} dove vi � una riga preliminare (comunque sempre dopo l'eventuale \emph{dichiarazione xml}) chiamata \emph{DTD} o \emph{Document Type Definition} tipo quella nel codice \ref{cod:dtd}.
\begin{code}[!h]
  \begin{blockcodexml}
    <!DOCTYPE html PUBLIC "-//W3C//DTD XHTML 1.0 Transitional//EN" "http://www.w3.org/TR/xhtml1/DTD/xhtml1-transitional.dtd">
  \end{blockcodexml}
  \caption{Esempio di DTD}\label{cod:dtd}
\end{code}
Se vi � la definizione dello schema allora oltre ad essere \emph{ben formato}, il documento deve essere anche \emph{valido}, ossia deve rispettare lo schema definito. Non approfondir� ulteriormente questa parte poich� nel mio programma mi occupo solo di controllare se un documento sia o meno \emph{ben formato}, ma ignoro eventuali definizioni di schema e la parte di validazione.

Nel mio programma, per motivi di abbreviazione e di chiarezza del codice pongo dei limiti ulteriori al set di possibili caratteri che � possibile usare, sia nel contenuto che nei nomi, ossia restringo il set alle lettere dell'alfabeto maiuscole e minuscole e ai numeri.
\section{Breve descrizione del programma}
Il programma prevede quattro funzionalit� di base, corrispondenti a quattro possibili \emph{goal} (notare la differenza nell'uso tra apici e virgolette):
\begin{itemize}
  \item\codeprolognb{controlla("stringa xml").} si ha un successo se la stringa indicata come parametro � \emph{ben formata}.
  \item\codeprolognb{controlla_file('nomeFile.xml').} si ha un successo se il file \emph{xml} indicato come parametro � \emph{ben formato}.
  \item\codeprolognb{completa("stringa xml").} si ha un successo se � possibile rendere \emph{ben formata} la stringa passata aggiungendoci codice prima e/o dopo (escludendo certi casi), inoltre viene visualizzato tale completamento.
  \item\codeprolognb{completa_file('nomeFile.xml').} si ha un successo se � possibile rendere \emph{ben formato} il file indicato aggiungendoci codice prima e/o dopo (escludendo certi casi), inoltre viene visualizzato tale completamento. 
\end{itemize}
tali \emph{goal} unificano con le prime quattro \emph{clausole} del programma, e come si pu� vedere la \emph{clausola} fulcro di tutte � \codeprolog{xml} o meglio \codeprolog{xml_} definita alla linea 43 del programma, visibile nella sezione \ref{sec:codice}, che con l'ausilio delle \emph{liste differenza} cerca di dividere il codice di input in tre parti: l'eventuale \emph{intestazione}, l'eventuale \emph{DTD}, e un singolo tag. Ignorer� il modello di esecuzione che porta all'unificazione dell'\emph{intestazione} e del \emph{DTD} in quanto � molto simile a quello del tag, quindi adesso passer� nel dettaglio a quest'ultimo.

Il singolo tag pu� unificare con una delle clausole definite alle righe: 80, 89, 97, una per un tag che contiene una lista di altri tags, una per un tag che contiene un contenuto, e una per un tag nella forma abbreviata vista prima nel codice \ref{cod:abbreviato}. Nel primo caso viene divisa la stringa in modo che al centro venga valutata una lista di tags (con la valutazione della clausola alla linea 103), circondata dall'apertura e chiusura del tag (alle linee 153 e 163) e da caratteri da ignorare come spazi, tabulatori, invio e anche i commenti (definito alla linea 109). In modo simile � definito anche il tag \emph{``foglia''} che contiene del testo invece di altri tags, al centro invece della valutazione di \codeprolog{tags} vi � quella di \codeprolog{contenuto} (alla linea 141) che ha successo se � una stringa costituita da possibili caratteri di contenuto.

Per quanto riguarda le clausole per l'\emph{apertura} e \emph{chiusura} dei \emph{tags}, si nota che vi � la variabile \codeprolog{Nome} che permette il controllo sulla coerenza di questo tra l'apertura e la chiusura. Entrambe le clausole tornano vero se le corrispondenti parti di tag sono correttamente formate (con una eventuale lista di attributi in apertura).

Si noti infine il modo in cui viene controllato se una stringa rispetta i possibili caratteri del possibile set, analizziamo solo il caso del nome di tag poich� pi� o meno gli altri casi (contenuto, attributi, etc\ldots) sono uguali. Nelle clausole di apertura e chiusura viene anche valutata (alle linee 155 e 165) la clausola \codeprolog{equal(A1-A2, Nome)} definita alla linea 190, che � vera se le stringhe \codeprolog{A1-A2} e \codeprolog{Nome} sono uguali e non nulle, inoltre per ogni carattere \codeprolog{H} viene valutata \codeprolog{carattereTag(H)} definito alla riga 238, che a sua volta torna vero se \codeprolog{H} appartiene alla stringa costituita dai possibili caratteri usabili per il nome.

Come si vede per esempio anche nella clausola \codeprolog{tag} vi sono dei punti in cui la ricorsione non � \emph{tail}, cio� durante l'esecuzione scende in profondit� nell'albero di valutazione e dopo deve valutare altri termini quando ritorna indietro, questo porta in certe occasioni durante la fase di \emph{completamento} ad una discesa su un ramo infinito dell'albero, e quindi ad una non risposta del programma. Un esempio di questo comportamento si pu� osservare nella chiamata \codeprolog{completa("</b></a>").}.
\section{Codice}\label{sec:codice}
\filecodeprolog{programma_per_relazione.pl}
\end{document}
